\documentclass[11pt,a4paper]{article}
\usepackage[utf8]{inputenc}
\usepackage[T1]{fontenc}
\usepackage[margin=2.5cm]{geometry}
\usepackage{parskip}
\pagestyle{empty}
\begin{document}
Dear Editor,

Please find enclosed our manuscript entitled "A conserved cysteine-histidine-glutamate metal site identifies DUF501 (Rv1025), an essential uncharacterised protein family of Mycobacterium tuberculosis, as a candidate metalloenzyme and drug target", which we submit for consideration as a Research Article in Antonie van Leeuwenhoek.

A quarter of a century after the Mycobacterium tuberculosis H37Rv genome was completed, a substantial part of its proteome remains functionally uncharacterised, and a subset of these dark proteins are essential for growth and highly vulnerable to partial knockdown, the profile sought for a drug target, yet they remain neglected because nothing is known about them. Our manuscript addresses one such protein, Rv1025, a 155-residue protein carrying the domain of unknown function DUF501 (Pfam PF04417), which is essential by transposon mutagenesis and strongly vulnerable by CRISPR interference. The entire DUF501 family, 4,370 proteins across bacteria, has no Gene Ontology term and no solved structure.

Combining population-scale sequence conservation, structural search against complete structural databases, genomic synteny across 63 Actinobacterial genomes, and AlphaFold3 complex and holo modelling, we show that DUF501 adopts a fold with no known structural relative, and we identify a near-invariant cysteine-histidine-glutamate triad that AlphaFold3 holo predictions consistently place a divalent metal on, with the coordination lost when the triad is mutated. Across the full Pfam alignment of the family, this triad proves strictly universal in bacteria, with no non-conservative substitution in any of the 2,228 sequences examined. We propose that DUF501 is a metal-binding protein and candidate metalloenzyme, the first functional hypothesis for this family, and that its conserved, essential metal pocket is a promising target for inhibition. As the study is fully computational, we present these findings as a hypothesis warranting experimental validation, not as proof of metal occupancy.

We believe this manuscript is well suited to Antonie van Leeuwenhoek because the journal's scope explicitly spans molecular biology and genetics through microbial pathogenesis to bioinformatics, and because the journal has recently published work of precisely this kind: the proteome-driven identification of candidate therapeutic targets in Orientia tsutsugamushi (2026), which likewise combined comparative genomic analysis across strains with druggability-based prioritisation of targets, without accompanying wet-lab validation. Our manuscript follows the same logic on a single, essential and vulnerable protein, and adds a structural argument that the earlier study did not attempt: a near-invariant metal-binding triad identified independently by two orthogonal predictors and calibrated for accessibility against validated catalytic zinc targets and a structural-zinc negative control. It also resolves, for the first time, a functional hypothesis for an entire bacterial domain family that currently carries no Gene Ontology term and no solved structure.

This manuscript has not been published previously and is not under consideration for publication elsewhere, in whole or in part. The author has read and approved the submitted version, has no competing interests to declare, and confirms the use of generative AI tools as part of the computational workflow, as disclosed in the Methods and Declarations sections of the manuscript.

Thank you for your consideration. I look forward to your response.

Sincerely,

Christophe Guyeux Femto-ST Institute, UMR 6174 CNRS, Universite Marie et Louis Pasteur, Besancon, France christophe.guyeux@univ-fcomte.fr
\end{document}
